\documentclass[aspectratio=169,10pt]{beamer}

% ---------- Theme ----------
\usetheme{Madrid}
\usecolortheme{default}
\setbeamertemplate{navigation symbols}{}

% ---------- Language and fonts ----------
\usepackage[T1]{fontenc}
\usepackage[utf8]{inputenc}
\usepackage[spanish,es-nodecimaldot]{babel}
\usepackage{lmodern}

% ---------- Utilities ----------
\usepackage{booktabs}
\usepackage{graphicx}
\usepackage{amsmath,amssymb}
\graphicspath{{figures/}}
\newcommand{\todo}[1]{\textcolor{red}{[TODO: #1]}}

% ---------- Metadata ----------
\title{Determinantes y forecasting de la variabilidad extrema del tipo de cambio}
\subtitle{Percentiles de retorno con ML y BART}
\author{Tu Nombre}
\institute{Tu Institucion}
\date{\today}

\begin{document}

\begin{frame}
  \titlepage
\end{frame}

\begin{frame}{Motivacion}
\begin{itemize}
    \item El riesgo cambiario extremo afecta inflacion, balance sheets y estabilidad financiera.
    \item Los modelos lineales suelen fallar en colas y durante cambios de regimen.
    \item ML y BART pueden capturar no linealidades e interacciones relevantes.
\end{itemize}
\end{frame}

\begin{frame}{Pregunta de investigacion}
\begin{block}{Objetivo}
Identificar determinantes y mejorar el pronostico de percentiles extremos de retornos del tipo de cambio.
\end{block}

\begin{itemize}
    \item Que variables importan en cola izquierda ($\tau=0.05$) y derecha ($\tau=0.95$)?
    \item Que metodo tiene mejor precision out-of-sample por horizonte?
\end{itemize}
\end{frame}

\begin{frame}{Datos y variable objetivo}
\begin{itemize}
    \item Frecuencia: \todo{diaria/semanal/mensual}.
    \item Muestra: \todo{fecha inicio - fecha fin}.
    \item Objetivo: $r_{t+h}=100(\log S_{t+h}-\log S_t)$.
    \item Quantiles de interes: $\tau\in\{0.01,0.05,0.10,0.90,0.95,0.99\}$.
\end{itemize}
\end{frame}

\begin{frame}{Variables explicativas}
\begin{columns}[T]
\column{0.5\textwidth}
\textbf{Bloque domestico}
\begin{itemize}
    \item tasas de interes
    \item inflacion esperada
    \item actividad economica
\end{itemize}

\column{0.5\textwidth}
\textbf{Bloque global/financiero}
\begin{itemize}
    \item VIX / estres financiero
    \item terminos de intercambio
    \item precio de commodities
\end{itemize}
\end{columns}
\end{frame}

\begin{frame}{Metodologia}
\begin{itemize}
    \item Benchmarks: ARX-Quantile Regression.
    \item ML: Random Forest, Gradient Boosting, XGBoost.
    \item No parametrico Bayesiano: BART para quantiles/colas.
    \item Evaluacion rolling/expanding window out-of-sample.
\end{itemize}
\end{frame}

\begin{frame}{Metrica principal}
\begin{equation*}
L_{\tau}(u)=u\left(\tau-\mathbb{I}\{u<0\}\right),\quad u=r_{t+h}-\hat{Q}_{\tau,t+h}
\end{equation*}

\begin{itemize}
    \item Menor pinball loss implica mejor pronostico por percentil.
    \item Se complementa con cobertura de eventos extremos.
\end{itemize}
\end{frame}

\begin{frame}{Resultados: desempeno por modelo}
\begin{table}
\centering
\begin{tabular}{lcccc}
\toprule
Modelo & $\tau=0.05$ & $\tau=0.10$ & $\tau=0.90$ & $\tau=0.95$ \\
\midrule
ARX-QR & -- & -- & -- & -- \\
RF-QR & -- & -- & -- & -- \\
GBM-QR & -- & -- & -- & -- \\
BART-QR & -- & -- & -- & -- \\
\bottomrule
\end{tabular}
\end{table}
\end{frame}

\begin{frame}{Resultados: episodios extremos}
\centering
\fbox{\parbox[c][0.45\textheight][c]{0.9\textwidth}{\centering
Inserte aqui la figura de quantiles pronosticados vs retornos observados}}

\vspace{0.3cm}
\small Quantiles pronosticados vs retornos observados.
\end{frame}

\begin{frame}{Interpretabilidad}
\begin{itemize}
    \item Importancia de variables por metodo (SHAP/permutation/BART inclusion frequency).
    \item Cambios de importancia entre cola izquierda y derecha.
    \item Implicaciones para monitoreo temprano de riesgo cambiario.
\end{itemize}
\end{frame}

\begin{frame}{Robustez}
\begin{itemize}
    \item Distintas ventanas y horizontes de pronostico.
    \item Submuestras por regimen de volatilidad.
    \item Verificacion de estabilidad temporal del ranking de modelos.
\end{itemize}
\end{frame}

\begin{frame}{Conclusiones}
\begin{itemize}
    \item \todo{Mensaje 1: principal determinante de cola extrema.}
    \item \todo{Mensaje 2: metodo ganador por horizonte/percentil.}
    \item \todo{Mensaje 3: uso practico para politica y gestion de riesgo.}
\end{itemize}
\end{frame}

\begin{frame}{Proximos pasos}
\begin{itemize}
    \item Afinar hiperparametros y estimar incertidumbre de forecast.
    \item Extender a panel de paises / multimoneda.
    \item Integrar nowcasting de variables de alta frecuencia.
\end{itemize}
\end{frame}

\begin{frame}
\centering
\Large Gracias
\end{frame}

\end{document}
